\section{EviRefine方法}\label{sec:method}
\subsection{问题定义与整体结构}
令$\X\in\mathbb R^{L\times C}$为多变量窗口，$\V$表示数据原本具有真值的位置，$\M\leq\V$表示人工遮挡后保留的可见位置。原生缺失不作为监督真值。以$\X_{\rm obs}$表示可见值，并将不可用位置置零。评价支撑与可见值保护定义为
\begin{equation}
 \widehat\X=\operatorname{PRISM}(\X_{\rm obs},\M).
 \label{eq:task}
\end{equation}

公式省略批次维度。推理使用输入窗口中的可见观测和符合资格的训练记忆。

\Model{}包含证据引导的粗恢复与观测条件化的残差精恢复两阶段（图\ref{fig:architecture}）。粗恢复生成训练集标准化尺度的显式估计$\X^{\rm c}$；精恢复将该估计作为固定输入并输出$\widehat\X$。粗恢复由TVCE、MPER和PEGR组成，OCRR承担第二阶段数值修正，共同恢复不完整输入窗口。

\subsection{时间--变量上下文编码（TVCE）}\label{sec:tvce}
TVCE结合局部时间动态与跨变量依赖，为各目标点建立上下文表示。
首先使用训练集统计进行变量标准化，再依据窗口内可见观测得到每个变量的均值$\mu_c$和标准差$\sigma_c$，必要时使用训练统计回退。数值变换前先施加掩码。共享滤波器的残差时间CNN逐变量处理可见数值、掩码、时间位置及到可见点的距离，并用全局水平信息补充局部表示。

变量签名注意力在完整时间特征签名之间交换信息。遵循T1的Channel--Head交互几何~\cite{t1}，每个特征通道对应一条变量注意力路径。模型包含两个时间残差块和两个变量注意力块，再对点级特征按重叠Patch池化：
\begin{equation}
 \begin{split}
  \F&=\operatorname{VSAttn}\bigl(\operatorname{CNN}(\X_{\rm obs},\M)\bigr),\\
  z_{j,c}&=\phi_{\rm patch}([P^x_{j,c};P^m_{j,c}])
    +\psi\bigl(\operatorname{Pool}(\F_{\mathcal P_j,c})\bigr).
 \end{split}\label{eq:context}
\end{equation}

其中$\operatorname{VSAttn}$表示变量签名注意力，$\mathcal P_j$表示第$j$个Patch，$P^x,P^m$是可见数值和掩码，$\psi$包含池化上下文的投影。固定$L=96$、Patch长度16、步长8、特征宽度64。基础重建头由$z_{j,c}$产生16维预测$b_{j,c}$。

\subsection{多来源Patch证据检索（MPER）}\label{sec:mper}
MPER利用来源可辨的互补片段增强目标上下文。
针对每个受损目标Patch，构建三类来源的候选集合$\mathcal C_{j,c}$。每个成员包含编码消息、质量特征、来源类型、非负来源内先验及有效标记。

\paragraph{时间Patch证据。}
候选来自当前窗口内同一变量的其他Patch。目标两侧分别按距离分档，最近档优先邻近位置，较远档依据学习到的匹配选择。没有可见支撑的Patch不具备候选资格。每侧五档，最多十个候选。

\paragraph{跨变量Patch证据。}
窗口相关的选择器将学习的变量关系与当前来源上下文结合，最多选择四个其他变量；每个来源提供时间偏移$-1,0,1$的片段，最多十二个候选。可见性及共同可见位置的一致性参与质量计算。候选编码中保留目标潜表示和来源上下文。

\paragraph{历史Patch证据。}
只由训练数据构建的记忆保存合法Patch成员及检索键。查询和检索辅助监督统一到训练集标准化坐标；局部归一化的查询在键编码前转换回该尺度。检索遵守来源变量与时间资格约束，从所选原型组返回最多四个成员，并将成员值转到目标工作尺度以构造消息。评价与精恢复训练时，发布的记忆快照和检索编码器保持固定。

统一证据库保留具体Patch成员、来源标记及支撑特征。质量特征与先验描述可观测支撑及兼容程度。当前结构在候选构建和逐点Query中均使用目标条件，其完整增量通过第\ref{sec:analysis}节的容量匹配上下文对照测量。

\subsection{逐点证据引导重建（PEGR）}\label{sec:pegr}
PEGR将Patch证据转化为各目标点对上下文基础估计的残差修正。
目标点利用点级上下文、Patch潜表示$z_{j,c}$、可见特征、片内位置、基础值及来源可用性构造Query $q_p$。候选消息经归一化后，加入质量投影和来源类型嵌入，再生成Key和Value。采用八个头，每头$d_h=8$维：
\begin{equation}
 \begin{split}
  a^{(h)}_{p,k}&=\frac{\exp(\ell^{(h)}_{p,k})}
   {1+\sum_{v\in\mathcal C_{j,c}}\exp(\ell^{(h)}_{p,v})},\\
  \ell^{(h)}_{p,k}&=\frac{\langle q^{(h)}_{p},K^{(h)}_k\rangle}{\sqrt{d_h}}
      +\log\pi_k+\eta_k,\\
  r_p&=\operatorname{Concat}_h\!\left(\sum_k a^{(h)}_{p,k}V^{(h)}_k\right),\\
  \delta_p&=\mathbf{1}_{\exists k\in\mathcal C_{j,c}:\eta_k=0}\,g_\theta([q_p;r_p]).
 \end{split}\label{eq:attention}
\end{equation}

$\pi_k$表示来源内先验；候选合法时$\eta_k=0$，否则为$-\infty$。分母中的常数1对应零logit、零Value的空候选，使证据不足时可以减弱读取量。残差头同时使用Query和拼接后的readout；没有合法候选时，修正为零。

将Patch估计变换到训练集标准化尺度，再通过置信度加权重叠融合合并，学习正的重叠权重：
\begin{equation}
 \begin{split}
 u_{j,p,c}&=\mu_c+\sigma_c b_{j,p,c}+s_c\delta_{j,p,c},\\
 s_c&=\sqrt{\sigma_c^2+0.1^2},\\
 w_{j,p,c}&=\exp\!\left(2\tanh\!\frac{h_{j,p,c}}2\right),\\
 X^{\rm c}_{t,c}&=\frac{\sum_{j:t\in\mathcal P_j}w_{j,t-jS,c}u_{j,t-jS,c}}
 {\sum_{j:t\in\mathcal P_j}w_{j,t-jS,c}}.
 \end{split}\label{eq:overlap}
\end{equation}

$s_c$中的数值下限避免可见窗口几乎不变时，残差尺度随之坍缩。$h_{j,p,c}$由置信度网络生成，输入包括点级上下文、可见特征、位置、基础/候选数值及来源可用性。这些权重决定各重叠Patch估计对融合结果的贡献。融合时排除补齐位置，并恢复原可见值。

\subsection{观测条件化残差精恢复（OCRR）}\label{sec:ocrr}
OCRR将时间上下文与真实可见观测提供的数值锚点结合，修正初始恢复轨迹。
第二阶段将固定粗估计与原可见观测结合。参考值$R_{t,c}$由同变量最近的左右可见点线性插值；只有单侧可见点时使用该点。支撑特征$U_{t,c}$随插值跨度或单侧距离衰减，不可用参考的支撑为零。

精恢复器将五类输入投影到64维，经过局部残差卷积及一层时间Transformer~\cite{attention}，产生两条修正：
\begin{equation}
 \begin{split}
 D&=U\Had(\R-\sg(\X^{\rm c})),\\
 Z&=\operatorname{TemporalEncoder}\bigl(
       [\sg(\X^{\rm c});\X_{\rm obs};\M;D;U]\bigr),\\
 \Delta_{\rm net}&=A_{\rm win}\Had2\tanh(W_n Z/2),\\
 G&=A_{\rm var}\Had\tanh(W_g Z),\\
 \widetilde\X&=\X^{\rm c}+\Delta_{\rm net}
      +G\Had(\R-\sg(\X^{\rm c})),\\
 \widehat\X&=\M\Had\X_{\rm obs}+(\mathbf 1-\M)\Had\widetilde\X.
 \end{split}\label{eq:refinement}
\end{equation}

该式的数值输入均位于训练集标准化坐标。$\sg$表示停止梯度，实际训练中整个粗模型冻结。激活标记为$A_{\rm win}=\mathbf{1}[\sum_{t,c}M_{t,c}>0]$和$A_{{\rm var},c}=\mathbf{1}[\sum_tM_{t,c}>0]$，按相应维度广播。第一条路径产生有界自由残差；第二条路径通过带符号门控$G$调节参考修正，只在该变量存在可见支撑时启用。整个窗口完全无观测时，网络修正关闭。最终掩码操作严格保留可见值。

\subsection{训练与选模}
先训练粗恢复，再固定其选定checkpoint训练精恢复器。在原本存在真值且被人工遮挡的位置，主要重建监督项为：
\begin{equation}
 \mathcal L_{\rm init}=\operatorname{MSE}_{\E}(\X^{\rm i},\X)+\lambda_k\mathcal L_{\rm key}+\lambda_l\mathcal L_{\rm level}.
 \label{eq:initial_objective}
\end{equation}
\begin{equation}
 \mathcal L_{\rm refine}=\frac{\|w\Had\E\Had(\widehat\X-\X)\|_1}{\|\E\|_1}.
 \label{eq:refinement_objective}
\end{equation}

训练结合上述重建项与辅助正则项；检索键监督将训练全局坐标中的掩码/完整键对齐与归一化对比监督结合。完整原生有效值只作为训练标签。两阶段均训练100轮，Adam学习率$10^{-3}$、权重衰减0、FP32、batch128。粗模型按固定验证集全窗口MSE选模，精恢复按全窗口MAE选模。精恢复训练期间，粗网络、记忆成员和检索编码器全部冻结。
